\documentstyle[%


righttag%


,11pt%


,amssymb%


,russian%


,izvru%


]{amsart}





\newcommand{\tl}[3]{\medskip\noindent{\bf#1}\ #2 \dotfill #3 \par} 


\pagestyle{empty}


\begin{document}


%\tableofcontents


\par


\vskip 2cm


\par


\bigskip


\bigskip


\bigskip


\par


\centerline{\Large СОДЕРЖАНИЕ}


\bigskip


\bigskip





\tl{Авсянкин О.Г., Карапетянц Н.К.}


{Об алгебре многомерных интегральных операторов\newline


с однородными ядрами с переменными коэффициентами}{3}


\tl{Андреев К.В.}


{О спинорном формализме при $n=6$}{11}


\tl{Габдулхаев Б.Г., Закиев М.И., Семенов И.П.}


{Оптимальные проекционные методы\newline


решения одного класса интегродифференциальных уравнений}{24}


\tl{Звягин В.Г., Дмитриенко В.Т., Кухарски З.}


{Топологическая характеристика\newline множества


решений фредгольмовых уравнений с $f$-компактно 


сужаемыми возмущениями\newline и ее приложения}{36}


\tl{Кельзон А.А.}


{Оценка погрешности тригонометрического


интерполирования функций\newline $m$-гармонической


ограниченной вариации}{49}


\tl{Шандра И.Г.}


{О конциркулярных тензорных полях


и геодезических отображениях\newline псевдоримановых пространств}{55}


\tl{Шишкин Г.И.}


{Сеточная аппроксимация волнового уравнения,


сингулярно возмущенного\newline по пространственной переменной}{67}


\medskip





\noindent{}Рефераты статей, депонированных в ВИНИТИ


через журнал ``Известия вузов.\newline Математика''.\dotfill{82}





\vfill


\noindent\raisebox{2pt}{\copyright} Казанский госудаpственный унивеpситет


\end{document}


